\section{Proofs}\label{app:proofs}
Throughout, $\|\cdot\|$ is the Euclidean norm and $\one$ the all-ones vector.

\subsection{Proof of Lemma~\ref{lem:quad}}
\emph{(a)} For one period, $R_i(\tau_i)=\tau_i(\alpha_i-\beta\tau_i)=\frac{\alpha_i^2}{4\beta}-\beta\big(\tau_i-\frac{\alpha_i}{2\beta}\big)^2$; sum over $i$.

\emph{(b)} Per period, with $x=(a-b\tau)/(1+bc)$ the inverse value function gives $\int_0^xv(z)\,dz-cx^2=\frac{ax}{b}-\frac{x^2}{2b}-cx^2=\frac{ax}{b}-\frac{1+2bc}{2b}x^2$. This is a concave quadratic in $x$ maximised at $x^*=a/(1+2bc)$, and the loss from any $x$ is $\frac{1+2bc}{2b}(x-x^*)^2$. The toll that implements $x^*$ solves $a-b\tau=(1+bc)a/(1+2bc)$, i.e.\ $\tau^W=ac/(1+2bc)$, and $x-x^*=-\frac{b}{1+bc}(\tau-\tau^W)$. Hence the loss is $\frac{1+2bc}{2b}\cdot\frac{b^2}{(1+bc)^2}(\tau-\tau^W)^2=\lambda(\tau-\tau^W)^2$. \qed

\subsection{Proof of Theorem~\ref{thm:known}}
Normalise $\alpha_2=1\le\alpha_1=r$ (the ratio is scale-invariant). Total revenue as a function of a common toll $\tau$ is, with $t=\beta\tau$,
\[R(\tau)=\frac{t}{\beta}\Big[(r-t)_++(1-t)_+\Big].\]
It is continuous and piecewise quadratic: for $t\le1$ it is $\frac1\beta t(r+1-2t)$, concave with maximiser $t_1=(r+1)/4$ and maximum $(r+1)^2/(8\beta)$; for $1\le t\le r$ it is $\frac1\beta t(r-t)$, concave with maximiser $t_2=r/2$ and maximum $r^2/(4\beta)$; for $t\ge r$ it is zero. Since $R$ is continuous, the global maximum is the larger of the maxima of the pieces restricted to their domains.

If $t_1\le1$ (i.e.\ $r\le3$), the first piece attains $(r+1)^2/(8\beta)$. If $t_1>1$ the first piece is increasing on its domain and its maximum is at $t=1$, which is also in the domain of the second piece, so it is dominated by the second piece's maximum over $[1,r]$. Likewise the second piece attains $r^2/(4\beta)$ iff $t_2\ge1$, i.e.\ $r\ge2$; otherwise its maximum over $[1,r]$ is at $t=1$, which is on the first piece.

\emph{Case $r\le2$.} Only the first piece has an interior maximiser, and the second piece is decreasing on $[1,r]$, so $R^S=(r+1)^2/(8\beta)$.

\emph{Case $2\le r\le3$.} Both maxima are attained, and $R^S=\max\{(r+1)^2/8,\,r^2/4\}/\beta$. Now $(r+1)^2\ge2r^2$ iff $r+1\ge\sqrt2r$ iff $r\le1+\sqrt2$.

\emph{Case $r\ge3$.} The first piece is increasing on $[0,1]$, so $R^S=r^2/(4\beta)$, and $r\ge3>1+\sqrt2$ is consistent with the claim.

Since $1+\sqrt2<3$ and $>2$, the three cases combine to $R^S=\frac{(r+1)^2}{8\beta}$ for $r\le1+\sqrt2$ and $\frac{r^2}{4\beta}$ beyond. The dynamic optimum is $R^D=(r^2+1)/(4\beta)$, so
\[\rho=\frac{(r+1)^2}{2(r^2+1)}\ (r\le1+\sqrt2),\qquad\rho=\frac{r^2}{r^2+1}\ (r\ge1+\sqrt2).\]
For the monotonicity: $\frac{d}{dr}\frac{(1+r)^2}{1+r^2}=\frac{2(1+r)(1-r)}{(1+r^2)^2}\le0$ for $r\ge1$, so the first branch decreases from $1$ at $r=1$ to $\rho(1+\sqrt2)$; explicitly $(2+\sqrt2)^2=6+4\sqrt2$ and $2(1+(1+\sqrt2)^2)=2(4+2\sqrt2)$, giving $\frac{6+4\sqrt2}{8+4\sqrt2}=\frac{2+\sqrt2}{4}=\frac12+\frac{1}{2\sqrt2}$. The second branch $r^2/(1+r^2)$ increases to $1$. Equality in $\rho=1$ requires $(\alpha_1-\alpha_2)^2=0$ in the first branch. \qed

\subsection{Proof of Theorem~\ref{thm:generalH}}
\emph{(a)} For $\tau\ge0$, $\tau(\alpha_i-\beta\tau)_+\ge\tau(\alpha_i-\beta\tau)$, so the true static revenue at the interior-optimal common toll $c=S/(2H\beta)$, $S=\sum\alpha_i$, is at least $\sum_ic(\alpha_i-\beta c)=S^2/(4H\beta)$. Thus $R^S\ge S^2/(4H\beta)$. The dynamic optimum $\tau_i=\alpha_i/(2\beta)$ is interior, $R^D=\sum_i\alpha_i^2/(4\beta)$. So $\rho\ge S^2/(H\sum_i\alpha_i^2)$.

For the second inequality use $(\alpha_i-m)(M-\alpha_i)\ge0$, i.e.\ $\alpha_i^2\le(m+M)\alpha_i-mM$. With $\mu=S/H$: $\sum_i\alpha_i^2\le H[(m+M)\mu-mM]$, so $\frac{S^2}{H\sum_i\alpha_i^2}\ge\frac{\mu^2}{(m+M)\mu-mM}$. Finally $\frac{\mu^2}{(m+M)\mu-mM}\ge\frac{4mM}{(m+M)^2}\iff(m+M)^2\mu^2-4mM(m+M)\mu+4m^2M^2\ge0\iff\big((m+M)\mu-2mM\big)^2\ge0$. The denominator is positive: since $\mu\ge m>0$, $(m+M)\mu-mM\ge(m+M)m-mM=m^2>0$.

\emph{(b)} Let $\alpha_i=\sqrt{H/i}$, $t=\beta\tau$. Then $R(\tau)=\frac t\beta\sum_{i\le k}(\alpha_i-t)$ where $k=k(t)$ is the number of $i$ with $\alpha_i>t$. Using $\sum_{i\le k}i^{-1/2}\le2\sqrt k$, $\sum_{i\le k}(\alpha_i-t)\le2\sqrt{Hk}-kt\le\max_{x\ge0}(2\sqrt{Hx}-xt)=H/t$ (maximum at $\sqrt x=\sqrt H/t$). Hence $R(\tau)\le H/\beta$ for every $\tau>0$, so $R^S\le H/\beta$. The dynamic revenue is $R^D=\frac1{4\beta}\sum_iH/i=\frac{HH_H}{4\beta}$. Therefore $\rho\le4/H_H$, and $H_H\sim\ln H\to\infty$. \qed

\subsection{Proof of Theorem~\ref{thm:est}}
Let $\tau^*=\tau^*(\theta)$, $\hat\tau=\tau^*(\hat\theta)=\tau^*+u$ (the map is affine, so $u$ is linear in the estimation error). Let $P=\frac1H\one\one^\top$ be the Euclidean projector onto the static subspace. The optimal static toll is $\bar\tau^*\one=P\tau^*$: for the revenue objective this is $\frac{S}{2H\beta}$, the mean of the $\alpha_i/(2\beta)$, and for any objective of the form of Lemma~\ref{lem:quad} the static optimum minimises $\|c\one-\tau^*\|^2$, i.e.\ is the mean. The static plug-in toll is $P\hat\tau=P\tau^*+Pu$ (again by averaging the estimated tolls).

(a) is Lemma~\ref{lem:quad}. For (b): the static plug-in's objective loss relative to the \emph{static} oracle is $\lambda\|P\hat\tau-P\tau^*\|^2=\lambda\|Pu\|^2=\lambda H\bar u^2$. For (c): $\|u\|^2=\|Pu\|^2+\|(I-P)u\|^2$ and $\|(I-P)u\|^2=\sum_i(u_i-\bar u)^2$. For (d): the static plug-in's loss relative to the \emph{dynamic} oracle is $\lambda\|P\tau^*+Pu-\tau^*\|^2=\lambda\|(I-P)\tau^*\|^2+\lambda\|Pu\|^2=G+\lambda H\bar u^2$ by orthogonality of $(I-P)\tau^*$ and $Pu$, while the dynamic plug-in's is $\lambda\|u\|^2$. The difference is $G+\lambda\|Pu\|^2-\lambda\|u\|^2=G-\lambda\|(I-P)u\|^2$. \qed

\subsection{Proof of Corollaries~\ref{cor:rev}--\ref{cor:cross} and Remark~\ref{rem:corr}}
For revenue, $\lambda=\beta$ and $u_i=\varepsilon_i/(2\beta)$, so $L^D=\beta\|u\|^2=(\varepsilon_1^2+\varepsilon_2^2)/(4\beta)$, $\bar u=(\varepsilon_1+\varepsilon_2)/(4\beta)$, $L^S=2\beta\bar u^2=(\varepsilon_1+\varepsilon_2)^2/(8\beta)$, and $L^D-L^S=\beta\frac{(u_1-u_2)^2}{2}=(\varepsilon_1-\varepsilon_2)^2/(8\beta)$. Also $G=\beta(\tau_1^*-\tau_2^*)^2/2=(\alpha_1-\alpha_2)^2/(8\beta)$. With i.i.d.\ errors of variance $\sigma_\alpha^2$, $\E L^D=\sigma_\alpha^2/(2\beta)$ and $\E L^S=\sigma_\alpha^2/(4\beta)$. For welfare, $u_i=c\varepsilon_i/(1+2bc)$ and $\lambda u_i^2=\frac{b(1+2bc)}{2(1+bc)^2}\cdot\frac{c^2\varepsilon_i^2}{(1+2bc)^2}=2\kappa\varepsilon_i^2$ with $\kappa=\frac{bc^2}{4(1+bc)^2(1+2bc)}$, which yields Corollary~\ref{cor:welf}.

For Corollary~\ref{cor:cross}, for $H=2$, $\sum_i(u_i-\bar u)^2=\frac12(u_1-u_2)^2$ and $G=\frac\lambda2(\tau_1^*-\tau_2^*)^2$, so the expected advantage is $\frac\lambda2[(\tau_1^*-\tau_2^*)^2-\E(u_1-u_2)^2]$. For revenue and i.i.d.\ errors, $\E(u_1-u_2)^2=\sigma_\alpha^2/(2\beta^2)$ and $(\tau_1^*-\tau_2^*)^2=(\alpha_1-\alpha_2)^2/(4\beta^2)$, giving the condition $\sigma_\alpha^2<(\alpha_1-\alpha_2)^2/2$.

For Remark~\ref{rem:corr}: with $\E u=\mu_u$ and $\Cov u=\Sigma_u$, $\E(u_1-u_2)^2=(\mu_{u,1}-\mu_{u,2})^2+\Var(u_1-u_2)$ and $\Var(u_1-u_2)=2\sigma_u^2(1-\varrho)$, so the static advantage $\frac\lambda2\E(u_1-u_2)^2$ equals $\lambda\sigma_u^2(1-\varrho)$ for centered errors. A common-mode bias leaves $u_1-u_2$ unchanged. \qed

\subsection{Proof of Theorem~\ref{thm:coupled}}
Revenue is $R(\tau)=\tau^\top\alpha-\tau^\top B\tau$, strictly concave with gradient $\alpha-2B\tau$ and Hessian $-2B$. Completing the square, $R(\tau)=\frac14\alpha^\top B^{-1}\alpha-(\tau-\tau^D)^\top B(\tau-\tau^D)$ with $\tau^D=\frac12B^{-1}\alpha$, which is (a). Restricting to $\tau=c\one$, $\|c\one-\tau^D\|_B^2$ is minimised at the $B$-orthogonal projection $P\tau^D=\frac{\one^\top B\tau^D}{\one^\top B\one}\one=\frac{\one^\top\alpha}{2\one^\top B\one}\one$, giving $c^S$; the value $R^S=R^D-\|(I-P)\tau^D\|_B^2=\frac14\alpha^\top B^{-1}\alpha-\|\tau^D\|_B^2+\|P\tau^D\|_B^2$, and $\|\tau^D\|_B^2=\frac14\alpha^\top B^{-1}\alpha$, so $R^S=\|P\tau^D\|_B^2=c^{S2}\one^\top B\one=(\one^\top\alpha)^2/(4\one^\top B\one)$, which is (b). The $B$-projector $P$ is self-adjoint in $\langle\cdot,\cdot\rangle_B$ and idempotent, so the Pythagorean identity $\|y\|_B^2=\|Py\|_B^2+\|(I-P)y\|_B^2$ holds.

(c) The plug-in dynamic toll is $\tau^D(\hat\alpha)=\tau^D+u$ with $u=\frac12B^{-1}\varepsilon$; its loss is $\|u\|_B^2=\frac14\varepsilon^\top B^{-1}BB^{-1}\varepsilon=\frac14\varepsilon^\top B^{-1}\varepsilon$. The plug-in static toll is the $B$-projection of the plug-in dynamic toll (by (b) applied to $\hat\alpha$, since $P$ is linear), i.e.\ $P\tau^D+Pu$, so its loss relative to the static oracle is $\|Pu\|_B^2=\frac{(\one^\top Bu)^2}{\one^\top B\one}=\frac{(\one^\top\varepsilon)^2}{4\one^\top B\one}$. The difference is $\|(I-P)u\|_B^2\ge0$, zero iff $u\in\mathrm{span}(\one)$ iff $B^{-1}\varepsilon\propto\one$ iff $\varepsilon\propto B\one$.

(d) As in Theorem~\ref{thm:est}: the static plug-in loses $\|(I-P)\tau^D\|_B^2+\|Pu\|_B^2$ relative to the dynamic oracle (the cross term vanishes by $B$-orthogonality), the dynamic plug-in loses $\|u\|_B^2$, and the advantage is $G-\|(I-P)u\|_B^2$ with $G=\|(I-P)\tau^D\|_B^2=R^D-R^S$. For mean-zero $\varepsilon$ with covariance $\Sigma$, $\E\|u\|_B^2=\frac14\tr(B^{-1}\Sigma)$ and $\E\|Pu\|_B^2=\frac{\one^\top\Sigma\one}{4\one^\top B\one}$. \qed

\subsection{Proof of Theorem~\ref{thm:bayes}}
Since $\tau^*$ is affine, $\E_\pi\tau^*(\theta)=\tau^*(\bar\theta)=:\bar\tau$. For any $\tau$, $\E_\pi\|\tau-\tau^*(\theta)\|^2=\|\tau-\bar\tau\|^2+\tr\Cov_\pi(\tau^*)$ (bias--variance decomposition). The posterior-expected objective is $\E_\pi\mathrm{Obj}^D(\theta)-\lambda\E_\pi\|\tau-\tau^*(\theta)\|^2$; maximising over $\tau\in\mathcal T$ minimises $\|\tau-\bar\tau\|^2$, whose unique solution is the Euclidean projection $\Pi_{\mathcal T}\bar\tau$ (part (i)). Substituting back gives (ii): the posterior expected loss against $\mathrm{Obj}^D(\theta)$ is $\lambda[\|(I-\Pi_{\mathcal T})\bar\tau\|^2+\tr\Cov_\pi\tau^*]$. Taking the dynamic class ($\Pi=I$) and the static class ($\Pi=P$) and subtracting yields $V=\lambda\|(I-P)\bar\tau\|^2\ge0$. \qed

\subsection{Proof of Theorem~\ref{thm:shrink}}
Write $\tau^*=\bar\tau^*\one+d$ with $\one^\top d=0$, and $\hat\tau=\tau^*+u$. Then $\hat\tau-\bar{\hat\tau}\one=d+(u-\bar u\one)$ and
\[\hat\tau(w)-\tau^*=\bar u\one+(w-1)d+w(u-\bar u\one).\]
The three vectors $\bar u\one$, $d$, $u-\bar u\one$ satisfy: $\bar u\one\perp d$ and $\bar u\one\perp(u-\bar u\one)$ pointwise, and $\E\langle d,u-\bar u\one\rangle=0$ because $\E u=0$. Hence $\E\|\hat\tau(w)-\tau^*\|^2=H\E\bar u^2+(1-w)^2D+w^2\E\|u-\bar u\one\|^2=s^2+(1-w)^2D+w^2(H-1)s^2$, which is (a). For (b), the quadratic in $w$ has minimiser $w^*=D/(D+(H-1)s^2)$ and minimum value $\frac{D(H-1)s^2}{D+(H-1)s^2}=\big[D^{-1}+((H-1)s^2)^{-1}\big]^{-1}$ (the parallel combination, strictly below both $D$ and $(H-1)s^2$). The comparison with the endpoints $w=0$ ($s^2+D$) and $w=1$ ($Hs^2$) follows.

(c) With $\tau_i^*=\mu+\delta_i$, $\delta_i\sim N(0,v)$ i.i.d.\ and a flat prior on $\mu$, the observations $\hat\tau_i=\mu+\delta_i+u_i$ are i.i.d.\ $N(\mu,v+s^2)$ given $\mu$. The posterior mean of $\tau_i^*$ is the standard empirical-Bayes shrinkage $\bar{\hat\tau}+\frac{v}{v+s^2}(\hat\tau_i-\bar{\hat\tau})$ (conditional on $\mu$ the posterior mean is $\mu+\frac v{v+s^2}(\hat\tau_i-\mu)$, and the flat prior replaces $\mu$ by $\bar{\hat\tau}$). With $\E D=(H-1)v$ the oracle weight is $\frac{(H-1)v}{(H-1)v+(H-1)s^2}=\frac v{v+s^2}$.

(d) The vector $z=\hat\tau-\bar{\hat\tau}\one$ is, in an orthonormal basis of $\one^\perp$, $N(d,s^2I_{H-1})$ and is independent of $\bar{\hat\tau}$. By James--Stein~\cite{jamesstein1961}, for dimension $q=H-1\ge3$ the estimator $(1-(q-2)s^2/\|z\|^2)z$ has strictly smaller risk than $z$ for every $d$, and the positive part (Baranchik) improves it further. Here $q-2=H-3$. The static component $\bar{\hat\tau}$ is unaffected (risk $s^2$). \qed

\subsection{Proof of Propositions~\ref{prop:ident-static}--\ref{prop:slope}}
\emph{Prop.~\ref{prop:ident-static}.} $Y_t+2\beta\tau=S+\xi_t$, so $\hat S=\frac1n\sum(Y_t+2\beta\tau)$ is unbiased with variance $\sigma^2/n$, and $\tau^S=S/(4\beta)$ is estimated with standard error $\sigma/(4\beta\sqrt n)$.

\emph{Prop.~\ref{prop:ident-split}.} By induction on $t$: given a toll history, the conditional law of $Y_t$ is $N$-type with mean $S-\beta(\tau_{1,t}+\tau_{2,t})$ and noise $\xi_t$, which depends on $\alpha$ only through $S$. If the toll rule is adaptive, $\tau_t$ is a (possibly randomised) function of $Y_1,\dots,Y_{t-1}$ with a fixed kernel, so the joint law of $(Y_1,\dots,Y_T)$ is the same under any two profiles with the same $S$, provided the interior condition holds in both so that the mean formula applies. Hence no statistic of the data separates them.

\emph{Example~\ref{ex:ident}.} Direct computation: $S=16$, $\tau^S=4$, $R^S=S^2/(8\beta)=32$; interiority for $\tau=4$: $6-4>0$ and $5-4>0$. $R^D=(100+36)/4=34$, $R^{D\prime}=(121+25)/4=36.5$; prizes over static $2$ and $4.5$. Interior for both profiles needs $\tau_1<10$ and $\tau_2<5$.

\emph{Remark~\ref{rem:corner-probe}.} At $(5.5,5.5)$: under $\alpha$, $x=(4.5,0.5)$, total $5$, revenue $27.5$; under $\alpha'$, $x=(5.5,0)$, total $5.5$, revenue $30.25$. The test between $N(5,\sigma^2)$ and $N(5.5,\sigma^2)$ from $n$ probe days has accuracy $\Phi(\Delta\mu\sqrt n/(2\sigma))$ with $\Delta\mu=0.5$. Per-day opportunity costs relative to $32$: $4.5$ and $1.75$ (using the static optimum of each profile, $32$ for both).

\emph{Prop.~\ref{prop:slope}.} Under the two probe levels, $\E\bar Y_\mp=S-2\beta\bar\tau\pm2\beta\delta$, so $\hat\beta=(\bar Y_--\bar Y_+)/(4\delta)$ is unbiased with variance $2(\sigma^2/n)/(16\delta^2)=\sigma^2/(8\delta^2n)$. Put $A=\frac12(\bar Y_-+\bar Y_+)$ with $\Var A=\sigma^2/(2n)$; $A$ and $\hat\beta$ are uncorrelated (sum and difference of two independent equal-variance variables). Then $\hat S=A+2\hat\beta\bar\tau$ has mean $S$ and variance $\frac{\sigma^2}{2n}+4\bar\tau^2\frac{\sigma^2}{8\delta^2n}=\frac{\sigma^2(\delta^2+\bar\tau^2)}{2\delta^2n}$. For $\hat\tau=\hat S/(4\hat\beta)=\frac{A}{4\hat\beta}+\frac{\bar\tau}2$, the delta method at $(\E A,\beta)=(\bar x,\beta)$, $\bar x=S-2\beta\bar\tau$, gives $\Var\hat\tau\approx\frac{\Var A}{16\beta^2}+\frac{\bar x^2}{16\beta^4}\Var\hat\beta=\frac{\sigma^2}{32\beta^2n}+\frac{\bar x^2\sigma^2}{128\beta^4\delta^2n}$, as stated. \qed

\subsection{Proof of Theorem~\ref{thm:review} and Corollaries~\ref{cor:daily},~\ref{cor:cost}}
For total $S$ and a toll $\tau$ the revenue is $\tau(S-2\beta\tau)$ and its loss against the static oracle $\tau=S/(4\beta)$ is $2\beta(\tau-S/(4\beta))^2$; with $\tau=\hat S/(4\beta)$ this is $(\hat S-S_t)^2/(8\beta)$. Index a cycle's days by $k=0,\dots,m-1$ and the previous cycle's days by $j=-m,\dots,-1$, with $S_{l+1}=S_l+\eta_l$. The estimate is $\hat S=\frac1m\sum_{j=-m}^{-1}S_j+\bar\xi$, with $\Var\bar\xi=\sigma^2/m$ independent of the $\eta$'s. For day $k$,
\[\hat S-S_k=\bar\xi-\sum_{l=-m}^{k-1}c_l\eta_l,\qquad c_l=\begin{cases}(l+m+1)/m,&-m\le l\le-1,\\1,&0\le l\le k-1.\end{cases}\]
(The increment $\eta_l$ lies between $S_j$ and $S_k$ iff $j\le l<k$; the number of such $j\in\{-m,\dots,-1\}$ is $l+m+1$ for $l\le-1$ and $m$ for $l\ge0$.) Therefore
\[\E(\hat S-S_k)^2=\frac{\sigma^2}m+\nu^2\Big[\sum_{i=1}^m\frac{i^2}{m^2}+k\Big]=\frac{\sigma^2}m+\nu^2\Big[\frac{(m+1)(2m+1)}{6m}+k\Big].\]
Averaging over $k=0,\dots,m-1$ (mean $k$ equal to $(m-1)/2$):
\[\frac{(m+1)(2m+1)}{6m}+\frac{m-1}2=\frac{2m^2+3m+1+3m^2-3m}{6m}=\frac{5m^2+1}{6m}.\]
Dividing by $8\beta$ gives $\bar L(m)$. Write $\bar L(m)=\frac1{8\beta}[\frac{5\nu^2}{6}m+\frac{\nu^2/6+\sigma^2}{m}]$, convex in $m>0$ with minimiser $m^{*2}=\frac{\nu^2/6+\sigma^2}{5\nu^2/6}=\frac{6\sigma^2+\nu^2}{5\nu^2}$ and minimum $\frac1{8\beta}\cdot2\sqrt{\frac{5\nu^2}6\big(\frac{\nu^2}6+\sigma^2\big)}=\frac{\nu\sqrt{5(6\sigma^2+\nu^2)}}{24\beta}$.

Corollary~\ref{cor:daily}: $\bar L(1)=\frac{\nu^2+\sigma^2}{8\beta}$, $\bar L(2)=\frac1{8\beta}[\frac{7\nu^2}4+\frac{\sigma^2}2]$; $\bar L(2)\ge\bar L(1)\iff\frac34\nu^2\ge\frac12\sigma^2\iff\sigma^2/\nu^2\le3/2$, and by convexity $m=1$ is the integer optimum iff this holds. The relaxed optimum equals $1$ iff $6\sigma^2+\nu^2=5\nu^2$, i.e.\ $\sigma^2/\nu^2=2/3$. As $m\to\infty$ the term $5\nu^2m/6\to\infty$. Corollary~\ref{cor:cost}: adding $C$ per revision, i.e.\ $C/m$ per day, replaces $\nu^2/6+\sigma^2$ by $\nu^2/6+\sigma^2+8\beta C$ in the numerator of the $1/m$ term, hence $m^{*2}=(6\sigma^2+\nu^2+48\beta C)/(5\nu^2)$. \qed

\subsection{Proof of Theorem~\ref{thm:kalman} and Proposition~\ref{prop:ewma}}
Let $\mathcal F_{-1}$ be the data up to the day before a revision day $0$, and let $\hat S^K$ be the steady-state Kalman prediction of $S_0$ from $\mathcal F_{-1}$; its error $\zeta=\hat S^K-S_0$ has variance $P=p+\nu^2$, where $p$ is the posterior variance satisfying $p=\frac{P\sigma^2}{P+\sigma^2}$, i.e.\ $p^2+\nu^2p-\nu^2\sigma^2=0$, whose positive root is as stated. Holding the toll $\hat S^K/(4\beta)$ for $k=0,\dots,m-1$ days gives error $\hat S^K-S_k=\zeta-\sum_{l<k}\eta_l$ with the $\eta_l$ ($l\ge0$) independent of $\mathcal F_{-1}$ hence of $\zeta$; its variance is $p+\nu^2+k\nu^2$, with mean over $k$ equal to $p+\nu^2+\frac{(m-1)\nu^2}{2}=p+\frac{(m+1)\nu^2}2$. This gives $\bar L_K(m)$ and (c).

(a) Let $\hat S$ be the cycle-mean estimate of $S_0$ from $\mathcal F_{-1}$. Its error $\hat S-S_0$ has second moment at least that of the best linear predictor, $P=p+\nu^2$ (the steady-state Kalman prediction is the $L^2$-projection of $S_0$ on the span of past observations, regardless of Gaussianity), and the identical orthogonality argument shows the cycle-mean policy's loss is $\E(\hat S-S_0)^2+\frac{(m-1)\nu^2}{2}\ge\bar L_K$. Equality would require the cycle-mean weights to equal the Kalman weights; the latter are geometric and positive at all lags, while the former vanish beyond $2m$ days and are constant on the window, so the inequality is strict. (Numerically, Table~\ref{tab:ksched}.) Adding $C/m$ to both sides preserves the strict inequality for every $m$, hence for the minimum over $m$.

(b) $\bar L_K(m)+C/m=\frac1{8\beta}[p+\frac{\nu^2}2]+\frac{\nu^2m}{16\beta}+\frac Cm$ is convex in $m$ with minimiser $m^{2}=16\beta C/\nu^2$ and minimum $\frac1{8\beta}(p+\frac{\nu^2}2)+2\sqrt{\frac{\nu^2C}{16\beta}}=\frac1{8\beta}(p+\frac{\nu^2}2)+\frac{\nu\sqrt C}{2\sqrt\beta}$. Over integers $m\ge1$ the minimiser is the better of the two integers around $\max\{1,m_K^*\}$; for $16\beta C<\nu^2$ revising daily is optimal.

\emph{Proposition~\ref{prop:ewma}.} Let $e_t=\hat S_t-S_t$. Then $\hat S_t-S_t=(1-g)(\hat S_{t-1}-S_{t-1}-\eta_{t-1})+g\xi_t$, so $e_t=(1-g)(e_{t-1}-\eta_{t-1})+g\xi_t$. The terms $e_{t-1}$, $\eta_{t-1}$, $\xi_t$ are independent, so the stationary variance $V$ solves $V=(1-g)^2(V+\nu^2)+g^2\sigma^2$, i.e.\ $V=\frac{(1-g)^2\nu^2+g^2\sigma^2}{g(2-g)}$. The toll for day $t+1$ uses $\hat S_t$ with error $e_t-\eta_t$ of variance $V+\nu^2$. For the random walk, the steady-state Kalman filter is $\hat S_t=\hat S_{t-1}+K(Z_t-\hat S_{t-1})$ with $K=P/(P+\sigma^2)$, an EWMA of gain $K^*=(p+\nu^2)/(p+\nu^2+\sigma^2)$; being the MMSE linear filter it minimises $V_g$ over $g$ with minimum $p$, which gives loss $(p+\nu^2)/(8\beta)$. \qed

\subsection{Proof of Propositions~\ref{prop:drift} and~\ref{prop:recovery}}
\emph{Prop.~\ref{prop:drift}.} With $S_j=S_0+\delta j$, $\hat S-S_k=\bar\xi+\delta(\overline{j}-k)$ where $\overline j=-(m+1)/2$, so the bias is $-\delta(k+\frac{m+1}2)$. Averaging its square over $k=0,\dots,m-1$: $\frac{(m-1)(2m-1)}6+\frac{(m+1)(m-1)}2+\frac{(m+1)^2}4=\frac{13m^2-1}{12}$ (common denominator $12$: $4m^2-6m+2+6m^2-6+3m^2+6m+3$). Adding the noise variance $\sigma^2/m$ and dividing by $8\beta$ gives the formula. Setting the derivative of $\frac{13\delta^2m^2}{12}+\frac{\sigma^2}m$ to zero gives $m^{*3}=\frac{6\sigma^2}{13\delta^2}$; at $m^*$, $\sigma^2/m^*=\frac{13\delta^2m^{*2}}6$, so the bracket is $\delta^2\frac{13m^{*2}-1+26m^{*2}}{12}=\frac{\delta^2(39m^{*2}-1)}{12}$.

\emph{Prop.~\ref{prop:recovery}.} Cycle $q$ comprises days $qm,\dots,qm+m-1$. If $j'=0$, the whole of cycle $q$ is post-jump, so the toll set at day $(q+1)m$ is the first based solely on post-jump data: lag $m$. If $j'>0$ cycle $q$ is mixed, cycle $q+1$ is pure post-jump, and the toll at day $(q+2)m$ is the first post-jump toll: lag $(q+2)m-J=2m-j'$. Averaging over $j'\in\{0,\dots,m-1\}$: $\frac1m[m+\sum_{j'=1}^{m-1}(2m-j')]=\frac1m[m+2m(m-1)-\frac{m(m-1)}2]=\frac{3m-1}2$. For estimate-once, the post-jump error is $\zeta-\Delta$, whose second moment is $\Delta^2+\E\zeta^2$ (for mean-zero $\zeta$), and the loss is that divided by $8\beta$. With $n_0$ initial days, $\E\zeta^2=\sigma^2/n_0$. \qed

\section{Additional notes on the experiments}\label{app:repro}
\paragraph{Exact optima.} Static optimal revenue with clipping is computed exactly by enumerating the finitely many breakpoints $\tau=\alpha_i/\beta$ and the stationary points of each piece, as in the proof of Theorem~\ref{thm:known}; no grid search is used. The dynamic optimum is the closed form. This makes E1 an exact check and not a statistical one.

\paragraph{Reporting standard.} A prediction $p$ is \emph{verified} against a simulated mean $\hat m$ with standard error $\mathrm{se}$ if $|p-\hat m|\le1.96\,\mathrm{se}$ or $|p-\hat m|\le0.005|p|$. Both clauses are needed: with $10^6$ or more draws the CI shrinks below the size of legitimate discretisation effects (initial transients, boundary of the cycle), and the relative clause is a pre-specified engineering tolerance, not tuned after seeing results. We report the number of failures.

\paragraph{Where simulations depart from the model.} (i) Losses use the true clipped revenue while plug-in rules use interior formulas (Section~\ref{sec:exp}). (ii) E3 uses stationary initial conditions at the first cycle boundary; the early transient is discarded. (iii) The A8 arm of the original protocol that studied queue-length observations was not implemented; only heteroskedastic noise proportional to flow was tested. (iv) E9 uses a Gaussian posterior truncated by omitting rows with $s/\mu>0.6$.

\paragraph{Compute.} Every experiment runs on a laptop CPU in minutes; seeds are fixed in the scripts, and tables are generated from a single JSON file so that every number in the paper has a one-to-one provenance.


